\documentclass[11pt,a4paper]{article}
\usepackage[T1]{fontenc}
\usepackage[utf8]{inputenc}
\usepackage[lithuanian]{babel}
\usepackage{lmodern}
\usepackage{microtype}
\setlength{\emergencystretch}{1.5em}
\usepackage{amsmath,amssymb,array,tabularx,enumitem,fancyhdr}
\usepackage[a4paper,margin=1.7cm,headheight=15pt]{geometry}
\setlength{\parindent}{0pt}
\setlength{\parskip}{5pt}
\setlist[itemize]{leftmargin=1.5em,itemsep=2pt,topsep=3pt}
\pagestyle{fancy}\fancyhf{}
\lhead{Kvadratinės lygtys: vertinimo instrukcija}\rhead{Mokytojui}
\cfoot{\thepage}
\newcommand{\uzd}[3]{\subsection*{#1 uždavinys \hfill #2 taškai}\textit{Sąlyga.} #3\par}
\newcommand{\kriterijai}{\textbf{Taškų skyrimas:}}
\newcommand{\balas}[2]{\item \textbf{#1 t.} #2}
\begin{document}
\begin{center}
{\LARGE\bfseries Vertinimo instrukcija}\par
{\large Poseidono variantas}\par
Varianto maksimumas -- \textbf{40 taškų}.
\end{center}
\textbf{Šaltinis.} Instrukcija parengta pagal paskutinėje žinutėje pateiktas uždavinių sąlygas, atsižvelgiant į vėlesnį Apolono 6 uždavinio patikslinimą. Išlaikyti jose nurodyti uždavinių taškai; jų vidinis paskirstymas sudarytas pagal vertinamus gebėjimus. Lygiaverčiai teisingi sprendimo būdai vertinami tuo pačiu maksimumu.

\section*{Bendrosios taisyklės}
\begin{itemize}
\item Taškai skiriami už parodytą matematinį turinį, ne už formulių ar eilučių kiekį. Sujungti, bet aiškiai pagrįsti žingsniai gauna visus atitinkamus taškus.
\item Vertinama atskirai kiekviename \textbf{savarankiškame vertinimo vienete}: visas uždavinys, o uždaviniuose su dalimis -- kiekviena a), b), c) dalis. Taigi antroji klaida vienoje dalyje nepanaikina teisingai išspręstos kitos dalies taškų. Tai šioje instrukcijoje pasirinkta klaidų taisyklės taikymo apimtis.
\item \textbf{Viena aritmetinė klaida.} Jei po vienos skaičiavimo klaidos mokinys nuosekliai ir matematiškai teisingai tęsia sprendimą su savo gautais skaičiais, atimamas tik \textbf{1 taškas} tame vertinimo vienete. Tos klaidos sukeltas neteisingas galutinis atsakymas papildomai nebaudžiamas.
\item \textbf{Antroji nepriklausoma aritmetinė klaida.} Jei tame pačiame vienete pasitaiko dar viena neištaisyta aritmetinė klaida, skiriami tik taškai už teisingai atliktus kriterijus \textbf{iki pirmosios aritmetinės klaidos}. Už pirmąją klaidą dar vienas taškas papildomai neatimamas. Už vėlesnę sprendimo dalį taškų neskiriama.
\item \textbf{Klaidos pasekmė nėra nauja klaida.} Pavyzdžiui, klaidinga šaknų suma, iš jos teisingai apskaičiuota kvadratų suma ir iš jos teisingai sudaryta lygtis -- viena klaida. Nauja klaida būtų neteisingai suskaičiuota paties mokinio gauta skaičių išraiška.
\item Aiškiai ištaisyta klaida nelaikoma neištaisyta klaida. Vertinamas galutinis, aiškiai pasirinktas sprendimo kelias. Jei pateikti keli prieštaringi atsakymai ir nė vienas nepasirinktas, galutinio atsakymo kriterijus neįvykdytas.
\item \textbf{Metodinė klaida} (neteisinga formulė, neteisėtas dalijimas, neteisingas keitinys, nepagrįstas sprendinių atmetimas) nevertinama kaip vien aritmetinė. Taškai už klaidingą metodą ir nuo jo priklausomas nepagrįstas išvadas neskiriami; nepriklausomi teisingi žingsniai gali būti įskaityti.
\item Vienos klaidos taisyklė nesuteikia taškų už neparodytus gebėjimus: jei klaida uždavinį supaprastino ir mokinys nebetikrino šakų ar sąlygų, neparodyti kriterijai neįskaitomi. Vertinimo minimumas -- 0 taškų.
\end{itemize}

\textbf{Tik atsakymas.} Be sprendimo galima skirti tik aiškiai atsakymui numatytus taškus. Išimtis -- tiesioginis skaidymo dauginamaisiais užrašas: teisinga pilna sandauga pati parodo prašytą rezultatą ir gauna visus tai daliai skirtus taškus.

\newpage
\section*{Kaip praktiškai taikyti klaidų taisyklę}
\textbf{Viena klaida:} nustatomas balas už parodytus teisingus metodus, vertinant vėlesnį skaičiavimą pagal mokinio gautus skaičius; iš šio balo atimamas 1 taškas. To paties klaidingo skaičiaus pasikartojimai ar klaidingas atsakymas antrą kartą nebaudžiami.

\textbf{Dvi klaidos:} grįžtama prie pirmosios klaidos vietos ir susumuojami tik iki jos teisingai įvykdyti kriterijai. Jei kriterijus reikalauja ir metodo, ir rezultato, o rezultatas iki pirmos klaidos dar negautas, už tą kriterijų taškas neskiriamas. Kai eilutėje numatyti 2 ar 3 taškai, jų vidinis paskirstymas žemiau nurodytas atskirai.

\subsection*{Pavyzdys: Apolono 5 uždavinys (5 t.)}
Teisingai reikia gauti $S=x_1+x_2=6$, $P=x_1x_2=4$, naują sumą $S^2-2P=28$ ir naują sandaugą $P^2=16$.
\begin{itemize}
\item Mokinys teisingai užrašo $S=6$, $P=4$, bet apskaičiuoja $36-8=30$. Toliau teisingai gauna naują sandaugą $16$ ir užrašo $y^2-30y+16=0$. \textbf{4/5 t.}: viena aritmetinė klaida, tolesnis metodas teisingas.
\item Tas pats mokinys papildomai užrašo $4^2=12$. \textbf{3/5 t.}: iki pirmos klaidos teisingai nustatė $S$ (1 t.), $P$ (1 t.) ir naujos sumos formulę $S^2-2P$ (1 t.). Taškai už vėlesnį darbą nebeskiriami.
\item Mokinys užrašo $x_1^2+x_2^2=(x_1+x_2)^2$. Tai \textbf{metodinė klaida}: praleistas $-2x_1x_2$ narys. Vieno taško aritmetinės klaidos lengvata čia netaikoma.
\end{itemize}

\subsection*{Kiti svarbūs sprendimo požymiai}
\begin{itemize}
\item Taškas už diskriminantą nėra privalomas sprendžiant skaidymu ar Vijeto teorema: jis pakeičiamas lygiaverčiu teisingu šaknų nustatymo žingsniu.
\item Už teisingas šaknis nereikalaujama dar ir įstatymo į pradinę lygtį, jei visas kelias buvo ekvivalentus ir apibrėžimo sritis nepažeista.
\item Sąlygose „neskaičiuodami sprendinių“ tikslas -- įvertinti darbą su suma ir sandauga. Vien tik aiškiai apskaičiuotų šaknų naudojimas nesuteikia taškų už prašytą šaknų neskaičiavimo metodą; galutinio rezultato taškai gali būti skiriami.
\item Galutinio atsakymo kriterijus suteikiamas vieną kartą. Jei du sprendiniai numatyti viename 1 taško kriterijuje, reikia abiejų; dalinis sprendinių radimas gali būti įskaitytas ankstesniame metodo kriterijuje.
\end{itemize}

\subsection*{Pažymio skyrimas}
\begin{center}\small
\begin{tabular}{|c|*{10}{c|}}\hline
Taškai&0--4&5--8&9--12&13--16&17--20&21--24&25--29&30--33&34--37&38--40\\\hline
Pažymys&1&2&3&4&5&6&7&8&9&10\\\hline
\end{tabular}\end{center}

\newpage
\section*{Poseidono variantas: 1--4 uždaviniai}
\uzd{1}{3}{$7x=4+3x^2$.}
\textbf{Sprendimas.} $3x^2-7x+4=0$, $(3x-4)(x-1)=0$. Atsakymas: $x=1$ arba $x=\frac43$.
\kriterijai\begin{itemize}
\balas{1}{Gauna $3x^2-7x+4=0$.}
\balas{1}{Teisingai skaido arba taiko kvadratinės lygties sprendimo metodą.}
\balas{1}{Pateikia abu sprendinius $1$ ir $4/3$.}
\end{itemize}
\uzd{2}{2}{Sudaryti redukuotą kvadratinę lygtį su šaknimis $\frac13,-5$.}
\textbf{Sprendimas.} $(x-\frac13)(x+5)=0\iff x^2+\frac{14}{3}x-\frac53=0$.
\kriterijai\begin{itemize}
\balas{1}{Teisingai nustato sumą $-14/3$ ir sandaugą $-5/3$ arba užrašo šaknų sandaugos formą.}
\balas{1}{Pateikia redukuotą lygtį $x^2+\frac{14}{3}x-\frac53=0$.}
\end{itemize}
\textit{Vertinimo pastaba.} $3x^2+14x-5=0$ yra lygiavertė, bet neredukuota lygtis: jei tik toks galutinis užrašas, daugiausia 1/2 t. Teisinga redukuota lygtis iš karto -- 2 t.
\uzd{3}{7}{a) $-2x^2+3x+5$; b) $x^3-5x^2-4x+20$; c) $x^3+2x^2-3$.}
\textbf{Sprendimas.} \textbf{a)} $-(2x-5)(x+1)$.\par\textbf{b)} $(x-5)(x^2-4)=(x-5)(x-2)(x+2)$.\par\textbf{c)} $x^3-1+2(x^2-1)=(x-1)(x^2+x+1)+2(x-1)(x+1)=(x-1)(x^2+3x+3)$. Pastarojo kvadratinio daugiklio $D=-3<0$.
\kriterijai\begin{itemize}
\balas{2}{a): 1 t. už tinkamas šaknis arba skaidymo pradžią; 1 t. už pilną sandaugą.}
\balas{2}{b): 1 t. už $(x-5)(x^2-4)$; 1 t. už pilną skaidymą.}
\balas{3}{c): 1 t. už šaknies $1$ atpažinimą ar tinkamą grupavimą; 1 t. už kvadratinio dalmens radimą; 1 t. už galutinę sandaugą $(x-1)(x^2+3x+3)$.}
\end{itemize}
\textit{Vertinimo pastaba.} Pilna teisinga sandauga iš karto gauna visus atitinkamos dalies taškus. c) nereikalaujama atskirai įrodinėti, kad kvadratinis daugiklis nebeskaidomas realiaisiais tiesiniais daugikliais.
\uzd{4}{4}{$P(x)=x^2+bx+9$ šaknų skirtumas $8$; rasti $b$.}
\textbf{Sprendimas.} Pagal Vijeto teoremą $S=-b$, $P=9$. Tada $(x_1-x_2)^2=S^2-4P=b^2-36=64$. Todėl $b^2=100$, $b=\pm10$. Abiem atvejais $D=64>0$.
\kriterijai\begin{itemize}
\balas{1}{Užrašo šaknų sumą $-b$ ir sandaugą $9$ arba naudoja formulę skirtumui per diskriminantą.}
\balas{1}{Gauna $(x_1-x_2)^2=b^2-36$.}
\balas{1}{Iš sąlygos gauna $b^2-36=64$, taigi $b^2=100$.}
\balas{1}{Pateikia abi reikšmes $b=-10$ ir $b=10$.}
\end{itemize}
\textit{Vertinimo pastaba.} Lygiavertis kelias: šaknų skirtumas yra $\sqrt D$, nes lygtis redukuota; todėl $\sqrt{b^2-36}=8$. Viena iš dviejų reikšmių nepakanka paskutiniam taškui.
\newpage
\section*{Poseidono variantas: 5--6 uždaviniai}
\uzd{5}{5}{$(x^2-5x)^2=2(x-2)(x-3)+12$.}
\textbf{Sprendimas.} $t=x^2-5x$, o dešinė pusė $2t+24$. Todėl $t^2-2t-24=0$, $(t-6)(t+4)=0$. Kai $t=6$, $x^2-5x-6=0$, $x=-1,6$. Kai $t=-4$, $x^2-5x+4=0$, $x=1,4$. Atsakymas: $x\in\{-1,1,4,6\}$.
\kriterijai\begin{itemize}
\balas{1}{Įveda keitinį ir gauna $t^2=2t+24$.}
\balas{1}{Randa abi reikšmes $t=6,-4$.}
\balas{1}{Teisingai išsprendžia šaką $t=6$: $x=-1,6$.}
\balas{1}{Teisingai išsprendžia šaką $t=-4$: $x=1,4$.}
\balas{1}{Pateikia visus keturis sprendinius.}
\end{itemize}
\textit{Vertinimo pastaba.} Jei išsprendžiama tik viena šaka, įskaitomi jos ir bendri pradžios kriterijai; pilno atsakymo taškas neskiriamas.
\uzd{6}{7}{$P(x)=-x^2+6x-3$, šaknys $m,n$; a) $m^2+n^2$; b) $1/m^3+1/n^3$; c) $m^2-n^2$, kai $m<n$.}
\textbf{Sprendimas.} $S=m+n=6$, $P=mn=3$.\par\textbf{a)} $S^2-2P=36-6=30$.\par\textbf{b)} $\dfrac{m^3+n^3}{(mn)^3}=\dfrac{S^3-3PS}{P^3}=\dfrac{216-54}{27}=6$.\par\textbf{c)} $(m-n)^2=S^2-4P=24$, todėl $m-n=-2\sqrt6$ ir $m^2-n^2=(m-n)S=-12\sqrt6$.
\kriterijai\begin{itemize}
\balas{2}{a): 1 t. už teisingą tapatybę su $S=6,P=3$; 1 t. už $30$.}
\balas{3}{b): 1 t. už bendrą vardiklį ir formulę kubų sumai; 1 t. už skaitiklį $162$ bei vardiklį $27$; 1 t. už rezultatą $6$.}
\balas{2}{c): 1 t. už neigiamo skirtumo nustatymą $m-n=-2\sqrt6$; 1 t. už $-12\sqrt6$.}
\end{itemize}
\textit{Vertinimo pastaba.} Kiekviena dalis yra atskiras klaidų taisyklės vienetas. Jei viena bendra aritmetinė klaida pradžioje paveikia kelias dalis, 1 t. už ją atimamas tik kartą visame uždavinyje; tos pačios klaidos pasekmės kitose dalyse nebaudžiamos. Pradinės klaidos priskyrimą fiksuoti vertinimo pastaboje.
\newpage
\section*{Poseidono variantas: 7--8 uždaviniai}
\uzd{7}{6}{$\dfrac{6-x}{x^2-4x+6}-\dfrac{4x^2-16x+24}{2x-12}=3$.}
\textbf{Sprendimas.} $Q=x^2-4x+6=(x-2)^2+2>0$, $x\ne6$. Imame $t=\frac{6-x}{Q}\ne0$. Kadangi antras narys su minusu lygus $\frac{2Q}{6-x}$, gauname $t+\frac2t=3$, $(t-1)(t-2)=0$.\par Kai $t=1$: $6-x=Q$, $x^2-3x=0$, $x=0,3$. Kai $t=2$: $6-x=2Q$, $2x^2-7x+6=0$, $x=\frac32,2$. Visi priklauso apibrėžimo sričiai. Atsakymas: $x\in\{0,\frac32,2,3\}$.
\kriterijai\begin{itemize}
\balas{1}{Nustato apibrėžimo sritį $x\ne6$; pagrindžia, kad $Q$ nelygus nuliui.}
\balas{1}{Pasirenka keitinį ir teisingai gauna $t+2/t=3$.}
\balas{1}{Gauna abi reikšmes $t=1,2$.}
\balas{1}{Išsprendžia pirmą šaką: $x=0,3$.}
\balas{1}{Išsprendžia antrą šaką: $x=3/2,2$.}
\balas{1}{Pateikia visus keturis tinkamus sprendinius, neįtraukdamas draudžiamų reikšmių.}
\end{itemize}
\textit{Vertinimo pastaba.} Galima be keitinio, padauginus iš vardiklių, gauti sandaugą $(Q-(6-x))(2Q-(6-x))=0$. Taškus priskirti atitinkamiems lygiaverčiams žingsniams.
\uzd{8}{6}{$x^2-mx+9-m=0$ turi realiąsias šaknis ir $x_1^2+x_2^2+x_1^2x_2^2=39$; rasti $m$ neskaičiuojant šaknų.}
\textbf{Sprendimas.} $S=m$, $P=9-m$. Sąlyga $S^2-2P+P^2=39$, todėl $m^2-2(9-m)+(9-m)^2=39$, $2m^2-16m+24=0$. Gauname $m=2,6$. Diskriminantas $D=m^2+4m-36$. Kai $m=2$, $D=-24<0$; kai $m=6$, $D=24>0$. Atsakymas: $m=6$.
\kriterijai\begin{itemize}
\balas{1}{Nustato sumą $S=m$ ir sandaugą $P=9-m$.}
\balas{1}{Perrašo sąlygą kaip $S^2-2P+P^2=39$.}
\balas{1}{Gauna $2m^2-16m+24=0$.}
\balas{1}{Randa abi kandidatines reikšmes $m=2,6$.}
\balas{1}{Patikrina realiųjų šaknų sąlygą: $D(2)=-24$, $D(6)=24$.}
\balas{1}{Atmeta $m=2$ ir pateikia vienintelį atsakymą $m=6$.}
\end{itemize}
\textit{Vertinimo pastaba.} Jei parašyta $m=2$ arba $m=6$, bet netikrintas realumas, maksimumas -- 4/6 t. Vien apskaičiuojant pačias šaknis -- neįskaitomi pirmų dviejų kriterijų taškai už prašytą metodą.
\end{document}
